\SourceLeaf{114}
\pdfbookmark[1]{Göttingische gelehrte Anzeigen. 1815 December 23.}{anzeige1815}
\begin{center}
\rule{0.65\linewidth}{0.4pt}\par
Göttingische gelehrte Anzeigen. 1815 December 23.\par
\rule{0.65\linewidth}{0.4pt}
\end{center}

Am 7. December wurde der Königl. Societät vom Herrn Prof. \textsc{Gauss} eine Vorlesung übergeben, überschrieben:

\begin{center}
\emph{Demonstratio nova altera theorematis, omnem functionem algebraicam rationalem integram unius variabilis in factores reales primi vel secundi gradus resolvi posse.}
\end{center}

Der hier ausgesprochene Lehrsatz, der wichtigste in der Theorie der Gleichungen, hat bekanntlich mehrere der ersten Geometer vielfältig beschäftigt. Die vornehmsten Versuche, einen vollkommen strengen Beweis davon zu geben, von \textsc{d'Alembert}, \textsc{Euler}, \textsc{Foncenex} und \textsc{Lagrange}, sind von dem Verf. gegenwärtiger Abhandlung bereits in einer vor sechszehn Jahren erschienenen Schrift zusammengestellt, und einer Prüfung unterworfen. So sehr man den Scharfsinn, welcher jene Arbeiten auszeichnet, anerkennen muss, so ist doch nicht zu leugnen, dass sie alle mehr oder weniger Lücken übrig lassen, wodurch die Beweiskraft zerstört wird, und wenn gleich durch \textsc{Lagrange}'s tiefeindringende Untersuchungen dem grössten Theile jener Mängel abgeholfen worden ist, so können doch auch die Bemühungen dieses grossen Geometers eben so wie die scharfsinnige Art, wie später \textsc{Laplace} diesen Gegenstand behandelt hat, gerade von dem Hauptvorwurfe, welcher alle jene versuchten Beweise trifft, nicht freigesprochen werden. Dieser besteht darin, dass man die Sache so genommen hat, als sei bloss die \emph{Form} der Wurzeln zu bestimmen, deren \emph{Existenz} man voraussetzte, ohne sie zu beweisen,
\SourceLeaf{115}
eine Schlussart, die gerade hier bloss illusorisch, und in der
That eine wahre \emph{petitio principii} ist. Alle die erwähnten Beweise
sind rein analytisch, auch der von \textsc{d'Alembert} nicht ausgenommen,
obgleich er in einer geometrischen Einkleidung erscheint, die aber
für die Sache selbst ganz gleichgültig ist; und man könnte daher sich
fast zu dem Schlusse verleiten lassen, als sei jene Voraussetzung,
deren Unzulässigkeit übrigens dem genannten Analysten selbst
entgangen ist, bei einer analytischen Behandlung gar nicht zu
vermeiden gewesen.

Der Verf. vorliegender Abhandlung hatte in der oben erwähnten Schrift
diesen Gegenstand auf eine ganz verschiedene Art behandelt, und
einen höchst einfachen neuen Beweis gegeben, welcher das
Eigenthümliche hat, dass er sich zum Theil auf die Geometrie der Lage
gründet, und übrigens in Ansehung der Strenge nichts zu wünschen übrig
zu lassen scheint. Zugleich hatte er aber schon damals erklärt, dass
er keinesweges für unmöglich halte, auf rein analytischem Wege zu
einem vollkommen strengen Beweise zu gelangen, und sich die
ausführliche Entwicklung seiner Ideen auf eine andere Zeit
vorbehalten. Andere Beschäftigungen hatten ihn bisher von diesem
Gegenstande abgezogen: die gegenwärtige Abhandlung ist bestimmt,
jene Zusage zu erfüllen.

Was nun diesen neuen Beweis selbst betrifft, so steht er dem erwähnten
frühern an Einfachheit und Kürze freilich sehr nach: allein dieser
Umstand liegt in der Beschaffenheit der subtilen Untersuchung selbst.
Immer bleibt es angenehm, neben einem höchst einfachen, aber zum
Theil auf Fremdartigem beruhenden Beweise des wichtigen Lehrsatzes,
noch einen zweiten zwar längern und künstlichern, aber für sich ganz
selbstständigen zu besitzen. Eine Darstellung der Hauptideen des neuen
Beweises, wenn sie nur einigermassen befriedigend ausfallen sollte,
würde übrigens bei weitem mehr Raum einnehmen, als der Zweck dieser
Blätter verstattet: wir müssen daher diejenigen, welche einen
Hauptreiz der Mathematik in der vollkommenen Klarheit ihrer Wahrheiten
und der höchsten Strenge ihrer Beweise finden, auf den bald
erscheinenden dritten Band der Commentationen der Societät verweisen.
\begin{center}\rule{0.28\linewidth}{0.4pt}\end{center}
